\documentclass[10pt,letterpaper]{article}
\usepackage[T1]{fontenc}
\usepackage[utf8]{inputenc}
\usepackage[spanish,es-tabla]{babel}
\usepackage{csquotes}
\usepackage{mathpazo}
\usepackage[scaled=0.9]{helvet}
\usepackage{courier}
\usepackage{calc}
\usepackage[top=2cm, bottom=2cm, left=2cm, right=2cm]{geometry}
\usepackage{xcolor}
\usepackage{booktabs}
\usepackage{longtable}
\usepackage{array}
\usepackage{ragged2e}
\usepackage{xurl}
\usepackage[strings]{underscore}
\usepackage{fancyhdr}
\usepackage{sectsty}
\usepackage{tocloft}
\usepackage[colorlinks=true,linkcolor=blue,urlcolor=blue,citecolor=blue,breaklinks=true]{hyperref}
\usepackage{enumitem}

\definecolor{uncpblue}{HTML}{1A237E}
\allsectionsfont{\color{uncpblue}}

\renewcommand{\cftsecleader}{\cftdotfill{\cftsecdotsep}}
\cftpagenumbersoff{section}
\cftpagenumbersoff{subsection}
\cftpagenumbersoff{subsubsection}
\renewcommand{\cftsecfont}{\normalfont}
\renewcommand{\cftsubsecfont}{\normalfont}
\renewcommand{\cftsubsubsecfont}{\normalfont}
\setlength{\parindent}{0pt}
\setlength{\parskip}{4pt}
\setlength{\tabcolsep}{3pt}
\renewcommand{\arraystretch}{0.85}

\pagestyle{fancy}
\fancyhf{}
\fancyhead[L]{\small\textcolor{uncpblue}{\textbf{SGSI-UNCP}}}
\fancyhead[R]{\small\textcolor{uncpblue}{P-SGSI-07}}
\fancyfoot[C]{\thepage}
\renewcommand{\headrulewidth}{0.4pt}
\renewcommand{\footrulewidth}{0.4pt}
\setlength{\headheight}{14pt}

\setlist{nosep,leftmargin=*,after=\vspace{2pt}}
\setlength{\emergencystretch}{2em}

\newcommand{\makecover}{
\begin{titlepage}
\centering
\vspace*{2cm}
{\Huge\bfseries\color{uncpblue} SGSI-UNCP\par}
\vspace{0.7cm}
{\LARGE Procedimiento de Eliminación Segura de Información y
Activos\par}
\vspace{0.5cm}
{\large Documento Base del Sistema de Gestión\\de Seguridad de la Información\par}
\vspace{1.5cm}
{\small Universidad Nacional del Centro del Perú\par}
\vfill
\end{titlepage}
}

\begin{document}

\makecover
\setcounter{tocdepth}{2}
\RaggedRight
\tableofcontents
\newpage

\RaggedRight
\section{Procedimiento de Eliminación Segura de Información y
Activos}\label{procedimiento-de-eliminaciuxf3n-segura-de-informaciuxf3n-y-activos}

\textbf{Organización:} Universidad Nacional del Centro del Perú (UNCP)
\textbf{Aprobado por:} Oficial de Seguridad y Confianza Digital
\textbf{Norma:} ISO/IEC 27001:2022 (Controles A.8.10, A.8.11, A.7.14)
\textbf{Referencia Técnica:} NIST SP 800-88 Rev.~1 (Guidelines for Media
Sanitization) \textbf{Alineamiento:} POL-SGSI-03 (Clasificación de la
Información), F-SGSI-05 (Acta de Eliminación Segura)

\begin{center}\rule{0.5\linewidth}{0.5pt}\end{center}

\subsection{Objetivo}\label{objetivo}

Establecer los métodos y controles para la eliminación segura de
información y activos de almacenamiento en la UNCP, garantizando que los
datos institucionales no puedan ser recuperados o reconstruidos una vez
finalizada su vida útil, mitigando el riesgo de fugas de información por
disposición inadecuada.

\subsection{Alcance}\label{alcance}

Este procedimiento aplica a toda información y activos de almacenamiento
de la UNCP, independientemente de su formato (físico o digital),
incluyendo:

\begin{itemize}
\tightlist
\item
  Documentos en papel (actas, expedientes, contratos, resoluciones).
\item
  Discos duros (HDD, SSD) de servidores, computadoras y laptops.
\item
  Medios extraíbles (USB, CD/DVD, tarjetas SD, cintas magnéticas).
\item
  Dispositivos móviles (smartphones, tablets) institucionales.
\item
  Equipos de red y almacenamiento (NAS, SAN, switches con almacenamiento
  interno).
\item
  Respaldos en cinta o discos externos.
\item
  Información en la nube (buckets, volúmenes, snapshots).
\end{itemize}

\subsection{Niveles de Eliminación
Segura}\label{niveles-de-eliminaciuxf3n-segura}

La UNCP adopta los tres niveles de sanitización definidos en
\textbf{NIST SP 800-88 Rev.~1}:

{\setlength{\RaggedRightRightskip}{0pt plus 5em}\small\sloppy \begin{longtable}[]{@{}
  >{\hspace{0pt}\RaggedRight\arraybackslash}p{(\linewidth - 6\tabcolsep) * \real{0.2500}}
  >{\hspace{0pt}\RaggedRight\arraybackslash}p{(\linewidth - 6\tabcolsep) * \real{0.2500}}
  >{\hspace{0pt}\RaggedRight\arraybackslash}p{(\linewidth - 6\tabcolsep) * \real{0.2500}}
  >{\hspace{0pt}\RaggedRight\arraybackslash}p{(\linewidth - 6\tabcolsep) * \real{0.2500}}@{}}
\toprule\noalign{}
\begin{minipage}[b]{\linewidth}\raggedright
Nivel
\end{minipage} & \begin{minipage}[b]{\linewidth}\raggedright
Método
\end{minipage} & \begin{minipage}[b]{\linewidth}\raggedright
Descripción
\end{minipage} & \begin{minipage}[b]{\linewidth}\raggedright
Cuándo Aplicar
\end{minipage} \\
\midrule\noalign{}
\endhead
\bottomrule\noalign{}
\endlastfoot
\textbf{Clear (Borrado Lógico)} & Sobrescritura con software
especializado (1 o más pasadas con datos aleatorios o ceros) & El medio
sigue siendo reutilizable internamente & Equipos que se reasignan a otra
área o usuario dentro de la UNCP \\
\textbf{Purge (Purgado)} & Desmagnetización (Degaussing) para HDD, o
borrado seguro criptográfico (Cryptographic Erase) para SSD & El medio
queda inutilizable para cualquier propósito posterior & Equipos que se
dan de baja definitiva para donación, venta o reciclaje \\
\textbf{Destroy (Destrucción Física)} & Trituración, pulverización,
incineración o perforación del medio & Destrucción irreversible del
medio físico & Medios dañados, clasificados como Confidenciales, o
cuando el riesgo de fuga de datos es alto \\
\end{longtable}\normalsize}

\subsection{Métodos Específicos por Tipo de
Medio}\label{muxe9todos-especuxedficos-por-tipo-de-medio}

\subsubsection{4.1 Medios Físicos
(Papel)}\label{medios-fuxedsicos-papel}

{\setlength{\RaggedRightRightskip}{0pt plus 5em}\small\sloppy \begin{longtable}[]{@{}
  >{\hspace{0pt}\RaggedRight\arraybackslash}p{(\linewidth - 4\tabcolsep) * \real{0.3333}}
  >{\hspace{0pt}\RaggedRight\arraybackslash}p{(\linewidth - 4\tabcolsep) * \real{0.3333}}
  >{\hspace{0pt}\RaggedRight\arraybackslash}p{(\linewidth - 4\tabcolsep) * \real{0.3333}}@{}}
\toprule\noalign{}
\begin{minipage}[b]{\linewidth}\raggedright
Clasificación
\end{minipage} & \begin{minipage}[b]{\linewidth}\raggedright
Método Exigido
\end{minipage} & \begin{minipage}[b]{\linewidth}\raggedright
Estándar
\end{minipage} \\
\midrule\noalign{}
\endhead
\bottomrule\noalign{}
\endlastfoot
\textbf{Pública} & Reciclaje convencional & N/A \\
\textbf{Uso Interno} & Trituración en partículas (corte cruzado, no
tiras) & DIN 66399 Nivel P-4 o superior \\
\textbf{Confidencial} & Trituración en partículas finas (corte cruzado,
partículas \textless{} 2x15 mm) & DIN 66399 Nivel P-5 o superior \\
\end{longtable}\normalsize}

\begin{itemize}
\tightlist
\item
  Las trituradoras deben estar ubicadas en áreas accesibles al personal
  autorizado y recibir mantenimiento periódico.
\item
  Para volúmenes masivos de documentos confidenciales, se puede
  contratar un servicio externo de destrucción certificada con emisión
  de certificado de destrucción.
\end{itemize}

\subsubsection{4.2 Discos Duros HDD}\label{discos-duros-hdd}

{\setlength{\RaggedRightRightskip}{0pt plus 5em}\small\sloppy \begin{longtable}[]{@{}
  >{\hspace{0pt}\RaggedRight\arraybackslash}p{(\linewidth - 4\tabcolsep) * \real{0.3333}}
  >{\hspace{0pt}\RaggedRight\arraybackslash}p{(\linewidth - 4\tabcolsep) * \real{0.3333}}
  >{\hspace{0pt}\RaggedRight\arraybackslash}p{(\linewidth - 4\tabcolsep) * \real{0.3333}}@{}}
\toprule\noalign{}
\begin{minipage}[b]{\linewidth}\raggedright
Destino
\end{minipage} & \begin{minipage}[b]{\linewidth}\raggedright
Método
\end{minipage} & \begin{minipage}[b]{\linewidth}\raggedright
Pasadas
\end{minipage} \\
\midrule\noalign{}
\endhead
\bottomrule\noalign{}
\endlastfoot
Reasignación interna (Clear) & Sobrescritura con software (DBAN,
Blanco·n·Cero o herramienta del fabricante) & 1 pasada con ceros \\
Baja/Donación (Purge) & Desmagnetización (Degausser) + verificación &
N/A \\
Baja/Reciclaje (Destroy) & Perforación mecánica (al menos 3 agujeros
atravesando los platos) o trituración industrial & N/A \\
\end{longtable}\normalsize}

\subsubsection{4.3 Unidades de Estado Sólido
SSD}\label{unidades-de-estado-suxf3lido-ssd}

{\setlength{\RaggedRightRightskip}{0pt plus 5em}\small\sloppy \begin{longtable}[]{@{}
  >{\hspace{0pt}\RaggedRight\arraybackslash}p{(\linewidth - 2\tabcolsep) * \real{0.5000}}
  >{\hspace{0pt}\RaggedRight\arraybackslash}p{(\linewidth - 2\tabcolsep) * \real{0.5000}}@{}}
\toprule\noalign{}
\begin{minipage}[b]{\linewidth}\raggedright
Destino
\end{minipage} & \begin{minipage}[b]{\linewidth}\raggedright
Método
\end{minipage} \\
\midrule\noalign{}
\endhead
\bottomrule\noalign{}
\endlastfoot
Reasignación interna & Cryptographic Erase (comando ATA SANITIZE o
herramienta del fabricante que invalide la clave de cifrado) \\
Baja/Donación & Cryptographic Erase + verificación; si el SSD no soporta
CE, destrucción física \\
Baja/Reciclaje & Destrucción física (trituración) \\
\end{longtable}\normalsize}

\textbf{Nota importante:} La sobrescritura tradicional no es efectiva en
SSD debido al desgaste de celdas y la reserva de bloques internos.
Siempre que sea posible, debe utilizarse Cryptographic Erase.

\subsubsection{4.4 Medios Extraíbles (USB, CD/DVD, Tarjetas
SD)}\label{medios-extrauxedbles-usb-cddvd-tarjetas-sd}

\begin{itemize}
\tightlist
\item
  Todos los medios extraíbles que hayan contenido información
  clasificada como \textbf{Confidencial} o \textbf{Uso Interno} deben
  ser destruidos físicamente al finalizar su vida útil.
\item
  Método aceptado: Trituración mecánica, perforación o incineración
  controlada.
\item
  No se acepta la reutilización de medios extraíbles que hayan contenido
  datos confidenciales.
\end{itemize}

\subsubsection{4.5 Dispositivos Móviles (Smartphones,
Tablets)}\label{dispositivos-muxf3viles-smartphones-tablets}

\begin{itemize}
\tightlist
\item
  \textbf{Android:} Restauración de fábrica + verificación de que el
  cifrado estaba activo (si no, destrucción física).
\item
  \textbf{iOS:} Borrado de contenido y configuración (activa el
  Cryptographic Erase del chip de cifrado).
\item
  \textbf{En caso de daño:} Destrucción física del dispositivo.
\end{itemize}

\subsubsection{4.6 Información en la
Nube}\label{informaciuxf3n-en-la-nube}

\begin{itemize}
\tightlist
\item
  Antes de dar de baja una suscripción cloud, la OTI debe:

  \begin{enumerate}
  \def\labelenumi{\arabic{enumi}.}
  \tightlist
  \item
    Verificar que todos los datos han sido migrados o respaldados según
    las políticas de retención.
  \item
    Eliminar todos los objetos (buckets, volúmenes, snapshots, bases de
    datos).
  \item
    Solicitar al proveedor la certificación de eliminación de los datos
    de sus servidores.
  \item
    Eliminar las cuentas de acceso asociadas.
  \end{enumerate}
\end{itemize}

\subsection{Verificación
Post-Eliminación}\label{verificaciuxf3n-post-eliminaciuxf3n}

Después de aplicar cualquier método de eliminación, se debe verificar:

{\setlength{\RaggedRightRightskip}{0pt plus 5em}\small\sloppy \begin{longtable}[]{@{}
  >{\hspace{0pt}\RaggedRight\arraybackslash}p{(\linewidth - 2\tabcolsep) * \real{0.5000}}
  >{\hspace{0pt}\RaggedRight\arraybackslash}p{(\linewidth - 2\tabcolsep) * \real{0.5000}}@{}}
\toprule\noalign{}
\begin{minipage}[b]{\linewidth}\raggedright
Medio
\end{minipage} & \begin{minipage}[b]{\linewidth}\raggedright
Método de Verificación
\end{minipage} \\
\midrule\noalign{}
\endhead
\bottomrule\noalign{}
\endlastfoot
Papel triturado & Inspección visual de que las partículas cumplen el
tamaño exigido \\
HDD sobrescrito & Lectura de verificación con la misma herramienta de
sobrescritura \\
HDD desmagnetizado & Comprobación con un medidor de campo o verificación
de que el disco no es reconocido por ningún sistema \\
SSD con Cryptographic Erase & Verificación de que la unidad aparece como
``no inicializada'' o sin particiones \\
Destrucción física & Inspección visual del daño irreversible \\
\end{longtable}\normalsize}

\subsection{Cadena de Custodia para
Eliminación}\label{cadena-de-custodia-para-eliminaciuxf3n}

Para activos clasificados como \textbf{Confidenciales} o de alto valor,
se debe mantener una cadena de custodia documentada:

\begin{enumerate}
\def\labelenumi{\arabic{enumi}.}
\tightlist
\item
  \textbf{Responsable} del activo solicita la baja mediante
  \textbf{F-SGSI-03}.
\item
  \textbf{OTI} recibe el activo, registra el estado y aplica el método
  de eliminación correspondiente.
\item
  \textbf{Testigo} (Oficial de Seguridad o delegado) presencia la
  eliminación y firma el \textbf{F-SGSI-05 (Acta de Eliminación
  Segura)}.
\item
  \textbf{Acta} se archiva en el expediente del activo por un período
  mínimo de \textbf{5 años}.
\end{enumerate}

\subsection{Proveedores Externos de
Eliminación}\label{proveedores-externos-de-eliminaciuxf3n}

Si se contrata a un tercero para la eliminación masiva de activos:

\begin{itemize}
\tightlist
\item
  El proveedor debe estar certificado (NAID AAA o equivalente).
\item
  Debe emitir un \textbf{Certificado de Destrucción} por lote con fecha,
  método, cantidad y tipo de medio.
\item
  La UNCP se reserva el derecho de realizar auditorías in situ al
  proceso del proveedor.
\item
  Se debe firmar un acuerdo de confidencialidad previo a la entrega de
  los activos.
\end{itemize}

\subsection{Responsabilidades}\label{responsabilidades}

{\setlength{\RaggedRightRightskip}{0pt plus 5em}\small\sloppy \begin{longtable}[]{@{}
  >{\hspace{0pt}\RaggedRight\arraybackslash}p{(\linewidth - 2\tabcolsep) * \real{0.5000}}
  >{\hspace{0pt}\RaggedRight\arraybackslash}p{(\linewidth - 2\tabcolsep) * \real{0.5000}}@{}}
\toprule\noalign{}
\begin{minipage}[b]{\linewidth}\raggedright
Rol
\end{minipage} & \begin{minipage}[b]{\linewidth}\raggedright
Responsabilidad
\end{minipage} \\
\midrule\noalign{}
\endhead
\bottomrule\noalign{}
\endlastfoot
\textbf{Custodio del activo} & Solicitar la baja y entregar el activo a
la OTI \\
\textbf{OTI} & Ejecutar o supervisar el método de eliminación; mantener
los registros \\
\textbf{Oficial de Seguridad} & Aprobar el método según la
clasificación; verificar la destrucción de activos Confidenciales \\
\textbf{Jefe de Área} & Autorizar la baja de activos de su unidad \\
\end{longtable}\normalsize}

\subsection{Documentos Relacionados}\label{documentos-relacionados}

{\setlength{\RaggedRightRightskip}{0pt plus 5em}\small\sloppy \begin{longtable}[]{@{}
  >{\hspace{0pt}\RaggedRight\arraybackslash}p{(\linewidth - 2\tabcolsep) * \real{0.5000}}
  >{\hspace{0pt}\RaggedRight\arraybackslash}p{(\linewidth - 2\tabcolsep) * \real{0.5000}}@{}}
\toprule\noalign{}
\begin{minipage}[b]{\linewidth}\raggedright
Código
\end{minipage} & \begin{minipage}[b]{\linewidth}\raggedright
Nombre
\end{minipage} \\
\midrule\noalign{}
\endhead
\bottomrule\noalign{}
\endlastfoot
POL-SGSI-03 & Política de Clasificación de Información y Respaldos \\
F-SGSI-03 & Formato de Baja de Usuario y Devolución de Activos \\
F-SGSI-05 & Acta de Eliminación Segura de Activos \\
POL-SGSI-05 & Política de Seguridad para Dispositivos Móviles del
Personal \\
D-SGSI-00 & Marco de Referencia Terminológico y Normativo \\
\end{longtable}\normalsize}

\begin{center}\rule{0.5\linewidth}{0.5pt}\end{center}

\textbf{Elaborado por:} Oficina de Tecnologías de la Información (OTI)
--- UNCP

\end{document}
